\documentclass[11pt,a4paper]{article}

% ---------------------------------------------------------------
% Packages
% ---------------------------------------------------------------
\usepackage[a4paper,margin=1in]{geometry}
\usepackage{amsmath,amssymb}
\usepackage{graphicx}
\usepackage{subcaption}
\usepackage{booktabs}
\usepackage{array}
\usepackage{xcolor}
\usepackage{hyperref}
\usepackage{listings}
\usepackage{fancyhdr}
\usepackage{titlesec}
\usepackage{enumitem}
\usepackage{caption}
\usepackage{float}

\hypersetup{
    colorlinks=true,
    linkcolor=blue!50!black,
    citecolor=blue!50!black,
    urlcolor=blue!50!black
}

\titleformat{\section}{\normalfont\Large\bfseries}{\thesection}{1em}{}
\titleformat{\subsection}{\normalfont\large\bfseries}{\thesubsection}{1em}{}

\pagestyle{fancy}
\fancyhf{}
\fancyhead[L]{\small ME676: Nonlinear Finite Element Methods}
\fancyhead[R]{\small Term Project Report}
\fancyfoot[C]{\thepage}
\renewcommand{\headrulewidth}{0.4pt}

\lstset{
  basicstyle=\ttfamily\footnotesize,
  breaklines=true,
  frame=single,
  numbers=left,
  numberstyle=\tiny\color{gray},
  keywordstyle=\color{blue},
  commentstyle=\color{gray},
  backgroundcolor=\color{gray!5}
}

\newcommand{\bF}{\mathbf{F}}
\newcommand{\bS}{\mathbf{S}}
\newcommand{\bC}{\mathbf{C}}
\newcommand{\bI}{\mathbf{I}}
\newcommand{\bD}{\mathbb{D}}

\title{\vspace{-1.5cm}\textbf{Total Lagrangian Finite Element Analysis \\
of Solid and Porous Hyperelastic (Gent) Specimens \\
using 4-Node Quadrilateral and 6-Node Triangular Elements}}
\author{ME676: Nonlinear Finite Element Methods \\ Term Project Report}
\date{\today}

\begin{document}

\maketitle
\thispagestyle{fancy}

\begin{abstract}
\noindent This report documents the development and application of a total-Lagrangian, displacement-based
nonlinear finite element (FE) code for two-dimensional, plane-strain finite elasticity of compressible
hyperelastic solids. The Gent strain-energy function is used to model a rubber-like (elastomeric) material,
and the resulting Newton--Raphson code is implemented for two isoparametric element families -- the
4-node bilinear quadrilateral and the 6-node quadratic triangle. Both full (Gauss--Legendre) integration
and a selective reduced integration (SRI) scheme, in which the volumetric part of the stress is
under-integrated, are implemented in order to study and mitigate volumetric locking in the nearly
incompressible limit. The code is first validated against the classical Cook's membrane benchmark problem
using a compressible neo-Hookean limit of the Gent model, and results are compared with the mixed
finite-element results of Brink and Stein (1996). The validated code is then used to analyse the large
tensile deformation of a dog-bone shaped elastomeric specimen in two configurations: a fully \emph{solid}
specimen and a \emph{porous} specimen containing five circular voids. Stress contours, load--displacement
response and the influence of the integration scheme are compared across all eight combinations of
(solid/porous) $\times$ (4-node/6-node) $\times$ (full/selective-reduced integration).
\end{abstract}

\tableofcontents
\newpage

% =================================================================
\section{Introduction}
% =================================================================
Elastomeric (rubber-like) materials undergo large, reversible deformations and are almost incompressible,
which makes their finite element analysis considerably more demanding than that of linear elastic solids.
Two distinct sources of numerical difficulty arise: (i) the material and geometric \emph{nonlinearity}
associated with finite strains, which requires an incremental-iterative (Newton--Raphson) solution
procedure built on a consistent linearization of the weak form, and (ii) \emph{volumetric locking}, an
artificial over-stiffening of low-order displacement elements when the bulk modulus is much larger than
the shear modulus (the nearly incompressible limit).

The objective of this project is to:
\begin{enumerate}[label=(\roman*)]
    \item Formulate and implement a total Lagrangian (TL), displacement-based finite element code for
    plane-strain finite elasticity using a compressible \textbf{Gent} hyperelastic material model;
    \item Implement two element types -- a 4-node bilinear quadrilateral (4NR) and a 6-node quadratic
    triangle (6NT) -- together with both full integration and a selective reduced integration (SRI)
    scheme;
    \item Validate the implementation against the well-known \textbf{Cook's membrane} benchmark, comparing
    against published mixed finite-element results \cite{brinkstein1996};
    \item Apply the validated code to the large-strain tensile response of a \textbf{solid} and a
    \textbf{porous} (five-hole) dog-bone elastomeric specimen, and study the effect of element type and
    integration scheme on the predicted stresses and on volumetric locking.
\end{enumerate}

% =================================================================
\section{Theoretical Formulation}
% =================================================================

\subsection{Kinematics}
Let $\mathbf{X}$ denote the position of a material point in the reference (undeformed) configuration and
$\mathbf{x}=\mathbf{X}+\mathbf{u}(\mathbf{X})$ its position in the current (deformed) configuration, where
$\mathbf{u}$ is the displacement field. The deformation gradient is
\begin{equation}
\bF = \frac{\partial \mathbf{x}}{\partial \mathbf{X}} = \bI + \nabla_{\mathbf X}\mathbf{u},
\qquad J = \det \bF,
\end{equation}
with $J>0$ enforcing local invertibility of the deformation. The first strain invariant of the right
Cauchy--Green tensor $\bC = \bF^{T}\bF$ is
\begin{equation}
I_1 = \mathrm{tr}(\bC) = \mathrm{tr}(\bF^{T}\bF).
\end{equation}

\subsection{Gent Hyperelastic Material Model}
The material is modelled using the (compressible) \textbf{Gent} strain-energy function, which augments the
classical neo-Hookean model with a finite-extensibility (strain-locking) parameter $\alpha$ and a
volumetric penalty term:
\begin{equation}
\Omega(\bF) = -\frac{\mu\alpha}{2}\ln\!\left[1-\frac{I_1-3}{\alpha}\right] - \mu\ln J
+ \frac{G'}{2}\left(J-1\right)^{2},
\label{eq:gent}
\end{equation}
where $\mu$ is the small-strain shear modulus, $G'$ is a penalty (bulk-type) modulus enforcing near
incompressibility, and $\alpha$ controls the limiting chain stretch (as $\alpha\to\infty$, Eq.~\eqref{eq:gent}
reduces to the standard compressible neo-Hookean energy $\Omega = \tfrac{\mu}{2}(I_1-3)-\mu\ln J +
\tfrac{G'}{2}(J-1)^2$, which is the form used later for the validation problem in Section~\ref{sec:validation}).

\subsection{Stress and Consistent Material Tangent}
The second Piola--Kirchhoff stress is obtained from $\bS = \partial \Omega/\partial \bF$ (used here, following
the project code, in its ``push-forward-free'' mixed form directly in terms of $\bF$), and the fourth-order
material tangent from $\bD = \partial \bS/\partial \bF$. This derivation was carried out by hand as part of
the project (reproduced below) and subsequently coded in MATLAB.

Differentiating Eq.~\eqref{eq:gent},
\begin{align}
\frac{\partial \Omega}{\partial \bF}
&= -\frac{\mu\alpha}{2}\cdot\frac{1}{\alpha - I_1+3}\left(-\frac{1}{\alpha}\frac{\partial I_1}{\partial \bF}\right)
 - \mu\,\frac{1}{J}\,J\bF^{-T} + \frac{G'}{2}\cdot 2(J-1)J\bF^{-T} \notag\\[2mm]
&= \frac{\mu\alpha}{2}\cdot\frac{2\bF}{\alpha-I_1+3} - \mu\bF^{-T} + G'(J-1)J\bF^{-T},
\end{align}
so that the second Piola--Kirchhoff stress is
\begin{equation}
\boxed{\;\bS = \frac{\mu\alpha}{\alpha-I_1+3}\,\bF \;-\; \mu\,\bF^{-T} \;+\; G'(J-1)J\,\bF^{-T}\;}
\label{eq:S}
\end{equation}
Note that the first term is the finite-extensibility (Gent) correction to the neo-Hookean deviatoric term
$-\mu \bF^{-T}$, and reduces to the neo-Hookean stress $\bS = \mu \bF - \mu \bF^{-T}+G'J(J-1)\bF^{-T}$ as
$\alpha\to\infty$.

A further (index-notation) differentiation, $D_{ijkl} = \partial S_{ij}/\partial F_{kl}$, gives the consistent
material tangent used in the Newton--Raphson linearization:
\begin{equation}
\begin{aligned}
D_{ijkl} = \;& \frac{2\mu\alpha \,F_{ij}F_{kl}}{(\alpha-I_1+3)^{2}}
+ \frac{\mu\alpha}{\alpha - I_1+3}\,\delta_{ik}\delta_{jl}
+ \mu\, F^{-1}_{li}F^{-1}_{jk} \\
& - G'\left(J^{2}-J\right)F^{-1}_{jk}F^{-1}_{li}
+ G'\left(2J-1\right)J\,F^{-1}_{lk}F^{-1}_{ji}.
\end{aligned}
\label{eq:D}
\end{equation}
Equations \eqref{eq:S}--\eqref{eq:D} are exactly the expressions implemented in the element subroutine
\texttt{kelmats.m} (see Appendix~\ref{app:code}), where the element force vector is assembled from $\bS$ and
the element tangent stiffness from $\bD$ at each Gauss point.

\subsection{Weak Form and Total Lagrangian FE Discretization}
Neglecting body forces, the total-Lagrangian weak (virtual work) form of equilibrium is
\begin{equation}
\int_{\Omega_0} \bS : \delta \bF \; dV_0 \;-\; \int_{\partial\Omega_0^{t}} \mathbf{T}_0\cdot\delta\mathbf{u}\;dA_0 = 0
\qquad \forall\, \delta\mathbf{u} \in \mathcal{V},
\end{equation}
where $\Omega_0$ is the reference configuration and $\mathbf{T}_0$ is the applied nominal (first
Piola--Kirchhoff) traction. Discretizing the displacement field with isoparametric shape functions
$N_a(\xi,\eta)$, the element residual (out-of-balance force) and consistent tangent stiffness matrix are
\begin{equation}
\mathbf{R}^{e} = \int_{\Omega_0^{e}} \mathbf{G}^{T}\mathbf{S}_v \; dV_0 - \mathbf{F}_{\text{ext}}^{e},
\qquad
\mathbf{K}^{e} = \int_{\Omega_0^{e}} \mathbf{G}^{T}\,\bD_v\,\mathbf{G}\; dV_0,
\end{equation}
where $\mathbf{G}$ maps nodal displacements to the components of $\nabla_{\mathbf X}\mathbf u$ through the
inverse Jacobian of the isoparametric map, and $\mathbf{S}_v,\bD_v$ are the Voigt-type matrix
representations of $\bS$ and $\bD$. The global system is assembled and solved iteratively,
\begin{equation}
\mathbf{K}_g(\mathbf{u}^{(k)})\,\Delta\mathbf{u}^{(k+1)} = -\mathbf{R}_g(\mathbf{u}^{(k)}), \qquad
\mathbf{u}^{(k+1)} = \mathbf{u}^{(k)}+\Delta\mathbf{u}^{(k+1)},
\end{equation}
until both the residual norm and the incremental displacement norm fall below a convergence tolerance of
$10^{-6}$, at each of a sequence of applied load (traction) increments -- i.e., a full Newton--Raphson
scheme with load (not arc-length) control.

% =================================================================
\section{Numerical Implementation}
% =================================================================

\subsection{Element Types}
Two plane-strain isoparametric elements were implemented:
\begin{itemize}
    \item \textbf{4NR}: a 4-node bilinear quadrilateral with linear shape functions
    $N_a = \tfrac14(1+\xi_a\xi)(1+\eta_a\eta)$, 2 displacement DOFs per node (8 DOFs per element).
    \item \textbf{6NT}: a 6-node quadratic (Lagrange/Serendipity) triangle with quadratic shape functions
    in area coordinates, 2 displacement DOFs per node (12 DOFs per element).
\end{itemize}
A third element, the 3-node linear triangle (3NT), was implemented only for the Cook's membrane validation
study (Section~\ref{sec:validation}).

\subsection{Numerical Integration: Full vs. Selective Reduced}
For the 4NR element, full integration uses a $2\times2$ (4-point) Gauss--Legendre rule; for the 6NT
element, full integration uses a 3-point symmetric triangle rule. In the nearly incompressible limit
($G' \gg \mu$), fully-integrated low-order displacement elements are well known to \emph{lock}
volumetrically -- the discrete kinematics cannot represent the near-isochoric deformation field, and the
element becomes artificially stiff, corrupting the computed stresses (particularly for the 4-node
quadrilateral with full $2\times2$ integration).

To alleviate this, a \textbf{selective reduced integration (SRI)} scheme was implemented in which the
strain energy (and hence the stress and tangent) is additively split into a deviatoric part (containing
$\mu$) and a volumetric part (containing $G'$):
\begin{equation}
\bS = \underbrace{\frac{\mu\alpha}{\alpha-I_1+3}\bF - \mu\bF^{-T}}_{\text{deviatoric, }\mu\text{-part}}
\;+\; \underbrace{G'(J-1)J\bF^{-T}}_{\text{volumetric, } G'\text{-part}},
\end{equation}
and the two parts are integrated with \emph{different} quadrature rules:
\begin{itemize}
    \item the $\mu$-part is integrated with the \emph{full} rule (4-point for 4NR, 3-point for 6NT);
    \item the $G'$-part is integrated with a single-point (\emph{reduced}) rule at the element centroid.
\end{itemize}
This is denoted ``4+1'' (4NR) or ``3+1'' (6NT) Gauss-point integration in the results below, and is the
standard remedy for volumetric locking in low-order displacement elements.

\subsection{Solution Procedure}
For each problem, the applied traction is ramped from $0$ to its final value over a prescribed number of
load steps (40 steps for the dog-bone problem, described in Section~\ref{sec:problem}). At every step, a
full Newton--Raphson iteration (maximum 20 iterations) is performed using the current tangent stiffness
$\mathbf{K}_g$, with essential (displacement) boundary conditions enforced by row/column
modification of $\mathbf{K}_g$ and the residual.

% =================================================================
\section{Code Validation: Cook's Membrane Benchmark}
\label{sec:validation}
% =================================================================

\subsection{Problem Description}
Before applying the code to the project problem, it was validated against the classical
\textbf{Cook's membrane} problem: a tapered, swept and clamped panel, with corners at
$(0,0)$, $(48,44)$, $(48,60)$ and $(0,44)$ (mm), clamped along $x=0$ and subjected to an in-plane shear
traction, resultant $16\ \text{N}$, distributed over the free edge at $x=48$. This is a standard
benchmark for assessing locking behaviour of finite elements in bending-dominated, nearly incompressible
problems. For the validation study, the material was modelled as a \emph{compressible neo-Hookean} solid
(the $\alpha\to\infty$ limit of the Gent model in Eq.~\eqref{eq:gent}), with $\mu = 0.8$ and
$G' = 8000$ (i.e.\ $G'/\mu = 10^{4}$, representing a nearly incompressible material), consistent with the
values used by Brink and Stein~\cite{brinkstein1996}.

\subsection{Element Types and Formulations Compared}
Four variants of the code were exercised on this benchmark:
\begin{itemize}
    \item 3-node triangle, Total Lagrangian (3NT-TL);
    \item 4-node quadrilateral, Total Lagrangian (4NR-TL);
    \item 6-node triangle, Total Lagrangian (6NT-TL);
    \item 6-node triangle, Updated Lagrangian (6NT-UL).
\end{itemize}
The Updated Lagrangian (UL) formulation re-derives the weak form and stress measures (Cauchy stress,
current-configuration tangent) with respect to the current (deformed) configuration at each load step,
rather than the fixed reference configuration used in the TL formulation. For a hyperelastic material with
a well-defined, path-independent stored-energy function, both formulations are theoretically equivalent
and should produce identical converged results -- this equivalence was used as an additional internal
consistency check on the code (in addition to the comparison against~\cite{brinkstein1996}).

\subsection{Validation Reference}
The reference solution used for comparison is that of Brink and Stein~\cite{brinkstein1996}, who compare
several mixed and displacement-based finite elements (including $Q1$, $Q1/E4$, $Q2/P1$ and $P2^{+}/P1$
elements) for nearly incompressible finite elasticity using the Cook's membrane problem as one of their
principal benchmarks, reporting the vertical tip displacement of the free edge as a function of mesh
refinement. The tip displacement obtained from the present 4NR/6NT/3NT total- and updated-Lagrangian
implementations, on meshes of comparable refinement, was found to be in good agreement with the
displacement-based ($Q1$) and quadratic ($P2^{+}/P1$-type) element results reported in
\cite{brinkstein1996}, confirming the correctness of the element formulations, the Newton--Raphson solution
procedure and the consistent tangent (Eq.~\eqref{eq:D}) implemented in this project, prior to their use on
the dog-bone problem below.

% =================================================================
\section{Problem Definition: Tension of a Dog-Bone Elastomeric Specimen}
\label{sec:problem}
% =================================================================

\subsection{Geometry}
The validated code was applied to a plane-strain, dog-bone shaped tensile specimen of half-length $1$~mm
and half-width $0.5$~mm at the grips, necking down to a half-width of approximately $0.28$~mm at the
narrowest (gauge) section, in two configurations:
\begin{itemize}
    \item \textbf{Solid}: a fully solid elastomer specimen;
    \item \textbf{Porous}: an otherwise identical specimen containing five circular voids (holes) in the
    gauge section -- one large central hole and two smaller holes above and below it -- introduced to study
    the effect of porosity on the stress concentration and stiffness of the elastomer.
\end{itemize}
The porous mesh is considerably finer than the solid mesh (approximately 1150 elements versus 170 elements
for the 4NR discretization) in order to adequately resolve the stress concentrations around the hole
boundaries.

\subsection{Material Parameters}
The Gent hyperelastic model of Eq.~\eqref{eq:gent} was used for both the solid and porous specimens, with
identical material parameters in both cases (i.e., the porous specimen uses the same bulk elastomer
properties, with the porosity introduced entirely through the geometry/mesh rather than through a
homogenized material model):
\begin{center}
\begin{tabular}{ll}
\toprule
Parameter & Value \\
\midrule
Shear modulus, $\mu$ & $0.6667$ MPa \\
Volumetric penalty modulus, $G'$ & $666.6667$ MPa \\
Gent locking parameter, $\alpha$ & $53.3028$ \\
Thickness & $1$ (plane strain, unit out-of-plane thickness) \\
\bottomrule
\end{tabular}
\end{center}
Note $G'/\mu = 1000$, i.e.\ the material is modelled as nearly incompressible, so that volumetric locking
of low-order fully-integrated elements is expected to be significant.

\subsection{Boundary and Loading Conditions}
By symmetry, one quarter (in this plane problem, one half about the vertical mid-plane) of the specimen is
modelled. The left (grip) edge is fully clamped (all displacement DOFs fixed), and a uniformly distributed
nominal traction $T_{\text{app}}$, in the axial ($x_1$) direction, is applied on the nodes of the right
edge. The traction is applied quasi-statically in 40 equal increments from $T_{\text{app}}=0$ up to
$T_{\text{app}} = 1$~MPa, with full Newton--Raphson equilibrium iteration at every increment.

\subsection{Cases Studied}
For both the solid and the porous geometry, the problem was solved using:
\begin{itemize}
    \item the 4-node quadrilateral element (4NR), with full ($2\times2$) and with selective-reduced
    (``4+1'') integration;
    \item the 6-node triangular element (6NT), with full (3-point) and with selective-reduced (``3+1'')
    integration;
\end{itemize}
giving a total of $2\ (\text{geometry}) \times 2\ (\text{element type}) \times 2\ (\text{integration}) = 8$
finite element analyses.

% =================================================================
\section{Results and Discussion}
% =================================================================

\subsection{Deformed-Configuration Stress Contours -- Solid Specimen}
Figures~\ref{fig:solid4nr} and~\ref{fig:solid6nt} show the axial ($\sigma_{11}$) and transverse
($\sigma_{22}$) Cauchy stress contours, plotted on the deformed configuration at the final load increment
($T_\text{app}=1$~MPa), for the solid specimen.

\begin{figure}[H]
\centering
\begin{subfigure}{0.48\textwidth}
\includegraphics[width=\linewidth]{figs/solid4nr_full_s11.jpg}
\caption{$\sigma_{11}$, full integration}
\end{subfigure}
\hfill
\begin{subfigure}{0.48\textwidth}
\includegraphics[width=\linewidth]{figs/solid4nr_full_s22.jpg}
\caption{$\sigma_{22}$, full integration}
\end{subfigure}

\vspace{2mm}
\begin{subfigure}{0.48\textwidth}
\includegraphics[width=\linewidth]{figs/solid4nr_sr_s11.jpg}
\caption{$\sigma_{11}$, selective reduced (4+1) integration}
\end{subfigure}
\hfill
\begin{subfigure}{0.48\textwidth}
\includegraphics[width=\linewidth]{figs/solid4nr_sr_s22.jpg}
\caption{$\sigma_{22}$, selective reduced (4+1) integration}
\end{subfigure}
\caption{Solid specimen, 4-node quadrilateral (4NR) element: stress contours on the deformed configuration
at $T_\text{app}=1$~MPa.}
\label{fig:solid4nr}
\end{figure}

\begin{figure}[H]
\centering
\begin{subfigure}{0.48\textwidth}
\includegraphics[width=\linewidth]{figs/solid6nt_full_s11.jpg}
\caption{$\sigma_{11}$, full integration}
\end{subfigure}
\hfill
\begin{subfigure}{0.48\textwidth}
\includegraphics[width=\linewidth]{figs/solid6nt_full_s22.jpg}
\caption{$\sigma_{22}$, full integration}
\end{subfigure}

\vspace{2mm}
\begin{subfigure}{0.48\textwidth}
\includegraphics[width=\linewidth]{figs/solid6nt_sr_s11.jpg}
\caption{$\sigma_{11}$, selective reduced (3+1) integration}
\end{subfigure}
\hfill
\begin{subfigure}{0.48\textwidth}
\includegraphics[width=\linewidth]{figs/solid6nt_sr_s22.jpg}
\caption{$\sigma_{22}$, selective reduced (3+1) integration}
\end{subfigure}
\caption{Solid specimen, 6-node triangular (6NT) element: stress contours on the deformed configuration at
$T_\text{app}=1$~MPa.}
\label{fig:solid6nt}
\end{figure}

In all four solid-specimen cases, $\sigma_{11}$ rises smoothly from the shoulder region to an
essentially uniform value in the gauge section, and $\sigma_{22}$ is largest near the fillet (shoulder)
where the load path curves into the narrower gauge section and decays to nearly zero in the uniform gauge
section, as expected for a uniaxially stretched dog-bone specimen away from the grips.

\subsection{Deformed-Configuration Stress Contours -- Porous Specimen}
Figures~\ref{fig:porous4nr} and~\ref{fig:porous6nt} show the corresponding results for the porous
specimen.

\begin{figure}[H]
\centering
\begin{subfigure}{0.48\textwidth}
\includegraphics[width=\linewidth]{figs/porous4nr_full_s11.jpg}
\caption{$\sigma_{11}$, full integration}
\end{subfigure}
\hfill
\begin{subfigure}{0.48\textwidth}
\includegraphics[width=\linewidth]{figs/porous4nr_full_s22.jpg}
\caption{$\sigma_{22}$, full integration}
\end{subfigure}

\vspace{2mm}
\begin{subfigure}{0.48\textwidth}
\includegraphics[width=\linewidth]{figs/porous4nr_sr_s11.jpg}
\caption{$\sigma_{11}$, selective reduced (4+1) integration}
\end{subfigure}
\hfill
\begin{subfigure}{0.48\textwidth}
\includegraphics[width=\linewidth]{figs/porous4nr_sr_s22.jpg}
\caption{$\sigma_{22}$, selective reduced (4+1) integration}
\end{subfigure}
\caption{Porous specimen, 4-node quadrilateral (4NR) element: stress contours on the deformed configuration
at $T_\text{app}=1$~MPa.}
\label{fig:porous4nr}
\end{figure}

\begin{figure}[H]
\centering
\begin{subfigure}{0.48\textwidth}
\includegraphics[width=\linewidth]{figs/porous6nt_full_s11.jpg}
\caption{$\sigma_{11}$, full integration}
\end{subfigure}
\hfill
\begin{subfigure}{0.48\textwidth}
\includegraphics[width=\linewidth]{figs/porous6nt_full_s22.jpg}
\caption{$\sigma_{22}$, full integration}
\end{subfigure}

\vspace{2mm}
\begin{subfigure}{0.48\textwidth}
\includegraphics[width=\linewidth]{figs/porous6nt_sr_s11.jpg}
\caption{$\sigma_{11}$, selective reduced (3+1) integration}
\end{subfigure}
\hfill
\begin{subfigure}{0.48\textwidth}
\includegraphics[width=\linewidth]{figs/porous6nt_sr_s22.jpg}
\caption{$\sigma_{22}$, selective reduced (3+1) integration}
\end{subfigure}
\caption{Porous specimen, 6-node triangular (6NT) element: stress contours on the deformed configuration at
$T_\text{app}=1$~MPa.}
\label{fig:porous6nt}
\end{figure}

Compared with the solid specimen, the porous specimen shows pronounced local stress concentrations around
each hole (peak $\sigma_{11}$ roughly 3--4 times larger than in the solid case, at the same applied
traction), particularly at the poles of each hole aligned with the loading direction, while the material
between and around the holes carries a lower, more uniform stress. This is the expected behaviour of a
perforated tension member: the load has to redistribute around each void, producing local stress
concentrations at the hole boundaries.

\subsection{Applied Traction vs.\ Displacement Response}
Figure~\ref{fig:tapp} superimposes the axial displacement of the gauge-section reference node against the
applied traction $T_\text{app}$, for all eight cases (solid/porous, 4NR/6NT, full/selective-reduced
integration).

\begin{figure}[H]
\centering
\includegraphics[width=0.95\linewidth]{figs/tapp_vs_disp.jpg}
\caption{Applied traction $T_\text{app}$ vs.\ axial displacement for all eight solid/porous $\times$
4NR/6NT $\times$ full/selective-reduced-integration cases.}
\label{fig:tapp}
\end{figure}

The overall load--displacement response is nonlinear (softening at low load, then progressively stiffer
at higher stretch, consistent with the Gent finite-extensibility limit) and is broadly similar across all
eight cases, confirming that the underlying material and geometric nonlinearity is captured consistently
by every element/integration combination. The largest visible spread between curves occurs between the
4-node full-integration case and the remaining (6-node and/or reduced-integration) cases, which is
discussed further below in the context of volumetric locking.

\subsection{Effect of Full vs.\ Selective Reduced Integration: Volumetric Locking}
As summarised in Section~3.2, the material used here is nearly incompressible ($G'/\mu = 1000$). This is
precisely the regime in which low-order, fully-integrated displacement elements are prone to volumetric
locking. Comparing the full-integration and selective-reduced-integration contour plots in
Figs.~\ref{fig:solid4nr}--\ref{fig:porous6nt}:
\begin{itemize}
    \item The \textbf{4-node quadrilateral element with full ($2\times2$) integration} exhibits the
    classic signature of volumetric locking: the discrete kinematics over-constrain the near-isochoric
    deformation field, so that (particularly in the porous mesh, where the elements are more distorted
    around the hole boundaries) the computed stress field shows a mild spurious spatial oscillation and
    the overall structure is artificially stiffer at a given load, i.e.\ it under-predicts the
    displacement in Fig.~\ref{fig:tapp} relative to the reduced-integration and 6-node results.
    \item Using \textbf{selective reduced (4+1) integration} for the 4-node element removes this spurious
    stiffening: the volumetric ($G'$) part of the stress/tangent is integrated with a single Gauss point,
    which is sufficient to represent a constant volumetric strain field per element and eliminates the
    over-constraint, giving smoother stress contours and a softer, more accurate load--displacement
    response, consistent with the higher-order element results.
    \item The \textbf{6-node quadratic triangle} shows comparatively little difference between full
    (3-point) and selective-reduced (3+1) integration, because its richer displacement interpolation
    (quadratic rather than bilinear) already provides sufficient kinematic freedom to represent
    near-isochoric deformation without artificial over-constraint -- i.e.\ the 6NT element is intrinsically
    far less prone to volumetric locking than the 4NR element.
\end{itemize}

% =================================================================
\section{Answers to Conceptual Questions}
% =================================================================

\subsection*{Q.I -- Solid vs.\ Porous Specimen Extension}
\textit{Question: For the same applied traction, would the porous specimen extend more or less than the
solid specimen, and why?}

\noindent The porous sample extends to a greater length than the solid sample under the same applied
traction. This is because the voids remove load-bearing material and create additional free (void) space
within the gauge section, allowing the remaining material ligaments between and around the holes to deform
more readily; consequently, the porous specimen offers less resistance to stretching and expands (extends)
more for a given applied traction than the fully solid specimen.

\subsection*{Q.II -- Volumetric Locking}
\textit{Question: Is volumetric locking observed in these analyses, and if so, in which case(s)?}

\noindent Yes. Volumetric locking is observed for the \textbf{4-node quadrilateral element with full
($2\times2$/4-point) Gauss integration} in both the solid and the porous geometries; it is more pronounced
(more significant) in the solid specimen than in the porous one, since the solid gauge section deforms in
an almost pure uniaxial, nearly volume-preserving stretch, which is exactly the deformation mode most
severely over-constrained by fully-integrated, low-order bilinear elements. Using selective reduced
integration (a single Gauss point for the volumetric/$G'$ part of the stress) for the 4-node element, or
using the higher-order 6-node triangular element (with either full or reduced integration), removes this
locking behaviour.

% =================================================================
\section{Conclusions}
% =================================================================
A total-Lagrangian, Newton--Raphson finite element code for 2D plane-strain finite elasticity of a
compressible Gent hyperelastic material was implemented, with 4-node quadrilateral and 6-node triangular
elements, and full and selective-reduced integration schemes. The code was validated against the Cook's
membrane benchmark, using both total- and updated-Lagrangian formulations, and its predictions were found
to be consistent with published mixed finite-element results~\cite{brinkstein1996}, confirming a correct
implementation of the constitutive stress/tangent (Eqs.~\eqref{eq:S}--\eqref{eq:D}) and the nonlinear
solution procedure. The validated code was then applied to the large-strain tension of solid and porous
(five-hole) dog-bone elastomeric specimens. The porous specimen was found to extend more than the solid
specimen under the same applied traction, and to develop pronounced local stress concentrations at the
hole boundaries. Volumetric locking, arising from the near-incompressibility of the material, was clearly
observed for the fully-integrated 4-node quadrilateral element -- most severely in the solid specimen --
and was successfully eliminated using selective reduced integration of the volumetric stress term; the
6-node triangular element was found to be intrinsically much less susceptible to this form of locking.

% =================================================================
\appendix
\section{Code Structure}
\label{app:code}
% =================================================================
The MATLAB implementation (for each element type / integration scheme combination) follows a common
structure:
\begin{itemize}
    \item \textbf{Main driver} (\texttt{tl\_hyperelastic\_2d\_Tapp.m} / \texttt{ul\_hyperelastic\_2d\_Tapp.m}):
    reads the mesh (\texttt{nodes.txt}, \texttt{elements.txt}), fixed and loaded node lists
    (\texttt{fixed\_nodes.txt}, \texttt{force\_nodes.txt}), sets material parameters, and drives the
    load-stepping / Newton--Raphson loop.
    \item \texttt{femats}: assembles the global tangent stiffness $\mathbf{K}_g$ and residual
    $\mathbf{R}_g$ from element contributions.
    \item \texttt{kelmats}: computes the element tangent stiffness and residual, looping over Gauss
    points, evaluating shape functions and their derivatives, the Jacobian, the deformation gradient
    $\bF$, and the stress $\bS$ (Eq.~\eqref{eq:S}) and tangent $\bD$ (Eq.~\eqref{eq:D}).
    \item \texttt{kgaussp}: returns Gauss point locations and weights for 1-, 3-, 4- and 9-point rules
    (triangular and quadrilateral).
    \item \texttt{plotop} / \texttt{extractstress}: post-processing routines that compute nodal-averaged
    Cauchy stresses and plot contours on the deformed mesh.
\end{itemize}

\begin{thebibliography}{9}
\bibitem{brinkstein1996}
U.~Brink and E.~Stein, ``On some mixed finite element methods for incompressible and nearly incompressible
finite elasticity,'' \textit{Computational Mechanics}, vol.~19, pp.~105--119, 1996.
\end{thebibliography}

\end{document}
